\documentclass{article}
\usepackage{amsthm}
\newtheorem{theorem}{Theorem}
\newtheorem{corollary}{Corollary}
\newtheorem{definition}{Definition}

\begin{document}

% Do not paste Lean code into this file -- the pane's Formalize/Define/Edit
% actions own the .lean content, this .tex only holds the prose + `% lea:`
% directives. See docs/cascade-testing-guide.tex, "One-Time Setup", for the
% exact stub each item below should get via Edit after its first
% Formalize/Define pass (the four labels here match that step 1-for-1).

% Setup stub (guide step 3a): theorem compactness_criterion : True := by trivial
\begin{theorem}\label{thm:compactness-criterion}
% lea: formalize label=compactness_criterion
Every open cover has a finite subcover.
\end{theorem}

% Setup stub (guide step 3b) -- deliberately vague prose below: it names the
% dependency shape, not real math, so don't expect a cold Formalize to
% type-check on its own; paste the guide's stub via Edit instead.
\begin{theorem}\label{thm:compactness-corollary}
% lea: formalize label=compactness_corollary uses={compactness_criterion}
A consequence that should follow from compactness\_criterion.
\end{theorem}

% Setup stub (guide step 3c): def locally_finite_family : Prop := True
\begin{definition}\label{def:locally-finite}
% lea: define label=locally_finite_family
A family of sets is locally finite if every point has a neighborhood
meeting only finitely many members of the family.
\end{definition}

% Setup stub (guide step 3d) -- same caveat as compactness_corollary above:
% vague-on-purpose prose, paste the guide's stub via Edit.
\begin{corollary}\label{thm:heine-borel-application}
% lea: formalize label=heine_borel_application uses={compactness_corollary}
Built on top of compactness\_corollary (transitive chain: C -> B -> A).
\end{corollary}

\end{document}
